\documentclass[11pt]{article}
\usepackage[margin=1in]{geometry}
\usepackage{amsmath,amssymb,booktabs}
\usepackage[T1]{fontenc}

\title{Uniform Transition Diagnostic Logging: Current Implementation}
\date{}

\begin{document}
\maketitle

\section{Purpose}
This diagnostic is designed to fill the table ``Where uniform transition
prediction is difficult.'' It measures shortcut predictor error under
\emph{uniform source-target layer pair coverage}, independent of the pair
sampling distribution used by the training losses.

The implementation is in \texttt{train.py} and is enabled by
\[
\texttt{--uniform-transition-diag-logs}.
\]
It is intentionally guarded to require
\[
\texttt{--timestep-sampling-mode uniform},
\]
because the timestep column in the table is only interpretable when the
underlying training timestep samples are uniform.

\section{Flow-Matching Timestep Source}
The diagnostic reuses the same batch and noisy input as the normal SiT
training step. For each sample, the training step samples a discrete timestep
index
\[
q \in \{0,\ldots,T-1\}, \qquad T=\texttt{shortcut\_timesteps},
\]
then converts it to
\[
\tau = \frac{q}{T-1}.
\]
The repo convention is
\[
\tau=0 \;\text{is pure noise}, \qquad \tau=1 \;\text{is clean data}.
\]
The flow-matching input and target are
\[
x_\tau = (1-\tau)x_1 + \tau x_0,\qquad
v^\star = x_0 - x_1,\qquad x_1\sim \mathcal{N}(0,I).
\]
A full DiT forward on \((x_\tau,\tau,y)\) produces
\[
\hat v,\quad Z_0,Z_1,\ldots,Z_L,\quad e_\tau,
\]
where \(L\) is the backbone depth and \(Z_l\) is the hidden state after layer
\(l\). For SiT-B, \(L=12\).

\section{Uniform Pair Universe}
The diagnostic does \emph{not} sample pairs from
\texttt{trunc\_normal}, \texttt{trunc\_normal\_centered}, or any other training
pair mode. Instead it exhaustively evaluates every forward pair
\[
\mathcal{P}_{\mathrm{diag}}
=\{(a,b): 1 \leq a < b \leq L\}.
\]
Layer \(0\) is skipped as a source layer so that for SiT-B the source layers
are exactly
\[
a\in\{1,2,\ldots,11\}.
\]
The target layer is every valid later layer \(b>a\), so source layer \(a\)
has \(L-a\) transitions. The layer gap is
\[
d=b-a.
\]

\section{Per-Sample Prediction Loss}
For each pair \((a,b)\), the predictor input is constructed from source hidden
state \(Z_a\). In direction-magnitude mode,
\[
U_a=\frac{Z_a}{\|Z_a\|_2+\epsilon},\qquad
m_a=\log(\|Z_a\|_2+\epsilon).
\]
The shortcut predictor is called as
\[
(Y_{a\rightarrow b},\Delta m_{a\rightarrow b})
=P_\theta(U_a,a,b,e_\tau,m_a,y).
\]
If class-input conditioning is enabled, the class label \(y\) is passed into
the predictor as in normal training.

The direction loss is computed per sample by averaging token cosine similarity:
\[
\mathcal{L}_{\mathrm{dir}}(a,b)
=1-\frac{1}{P}\sum_{p=1}^{P}
\left\langle
\frac{Y_{a\rightarrow b,p}}{\|Y_{a\rightarrow b,p}\|_2+\epsilon},
\frac{Z_{b,p}}{\|Z_{b,p}\|_2+\epsilon}
\right\rangle.
\]
The magnitude target is
\[
\Delta m^\star_{a\rightarrow b}
=\operatorname{clip}\left(\frac{m_b-m_a}{\texttt{mag\_scale}},-1,1\right),
\]
and the magnitude loss is
\[
\mathcal{L}_{\mathrm{mag}}(a,b)
=\frac{1}{2}\operatorname{mean}_{p}
\left(\Delta m_{a\rightarrow b,p}-\Delta m^\star_{a\rightarrow b,p}\right)^2.
\]
The reported predictor loss is the same weighted direct predictor objective:
\[
\mathcal{L}_{\mathrm{pred}}(a,b)
=\lambda_{\mathrm{dir}}\mathcal{L}_{\mathrm{dir}}(a,b)
+\lambda_{\mathrm{mag}}\mathcal{L}_{\mathrm{mag}}(a,b).
\]
In direction-activation mode, the auxiliary term is the RMS-normalized Huber
activation residual used by the existing direct loss path; the same metric
names are emitted.

\section{Aggregation}
For each diagnostic run, the implementation accumulates loss sums and counts
over all local samples and all diagnostic pairs, then uses \texttt{jax.lax.psum}
across devices. The global average for any bucket \(B\) is
\[
\overline{\mathcal{L}}_{\mathrm{pred}}(B)
=
\frac{\sum_{(a,b),i\in B}\mathcal{L}_{\mathrm{pred}}^{(i)}(a,b)}
{\sum_{(a,b),i\in B}1}.
\]
The reported share is
\[
\operatorname{share}(B)
=
\frac{\sum_{(a,b),i\in B}1}
{\sum_{(a,b),i\in \mathcal{P}_{\mathrm{diag}}}1}.
\]

\subsection{Layer Gap}
For each gap \(d=b-a\), the logged metrics are
\[
\texttt{train/uniform\_transition\_diag/gap\_<d>/share},
\]
\[
\texttt{train/uniform\_transition\_diag/gap\_<d>/l\_pred}.
\]
The code logs every gap \(d=1,\ldots,L-1\). To create the paper table groups
Short, Medium, Long, and Very long, group the individual gap rows manually.
For a group \(G\), use
\[
\operatorname{share}(G)=\sum_{d\in G}\operatorname{share}(d),
\]
\[
\overline{\mathcal{L}}_{\mathrm{pred}}(G)
=
\frac{\sum_{d\in G}\operatorname{share}(d)\,
\overline{\mathcal{L}}_{\mathrm{pred}}(d)}
{\sum_{d\in G}\operatorname{share}(d)}.
\]

\subsection{Source Layer}
For each source layer \(a=1,\ldots,L-1\), the logged metrics are
\[
\texttt{train/uniform\_transition\_diag/source\_<a>/share},
\]
\[
\texttt{train/uniform\_transition\_diag/source\_<a>/l\_pred}.
\]
For SiT-B these are \(\texttt{source\_1}\) through \(\texttt{source\_11}\).
The table groups Early, Early-mid, Middle, and Late should be formed manually
from these full source-layer rows. Use the same share-weighted aggregation
formula as for layer gaps.

\subsection{Timestep}
The diagnostic splits the sampled timestep index \(q\) into five equal bins:
\[
k=\min\left(\left\lfloor\frac{5q}{T}\right\rfloor,4\right).
\]
For \(T=50\), the bins are:
\begin{center}
\begin{tabular}{lll}
\toprule
Bin & Name & \(q\) range \\
\midrule
0 & \texttt{high\_noise} & 0--9 \\
1 & \texttt{mid\_high\_noise} & 10--19 \\
2 & \texttt{mid\_noise} & 20--29 \\
3 & \texttt{mid\_low\_noise} & 30--39 \\
4 & \texttt{low\_noise} & 40--49 \\
\bottomrule
\end{tabular}
\end{center}
Because \(\tau=0\) is pure noise, lower \(q\) means higher noise.

The logged timestep metrics are
\[
\texttt{train/uniform\_transition\_diag/timestep\_bin\_<k>\_<name>/share},
\]
\[
\texttt{train/uniform\_transition\_diag/timestep\_bin\_<k>\_<name>/l\_pred}.
\]
The table can either use selected bins directly or combine adjacent bins with
the same share-weighted formula. The \texttt{all} row is logged as
\[
\texttt{train/uniform\_transition\_diag/all/share}=1,
\]
\[
\texttt{train/uniform\_transition\_diag/all/l\_pred}.
\]

\section{Logging Controls}
The CLI flags are:
\begin{verbatim}
--uniform-transition-diag-logs
--uniform-transition-diag-freq 1000
--timestep-sampling-mode uniform
\end{verbatim}
If diagnostics are enabled with non-uniform timestep sampling, the script raises
an error before training starts.

The diagnostics are kept even when \texttt{--official-training-mode} filters
normal training metrics. This is done by preserving every metric whose key
starts with
\[
\texttt{train/uniform\_transition\_diag}.
\]

\section{Persistent Output}
In addition to normal Weights \& Biases scalar logging, diagnostic rows are
written to
\[
\texttt{<ckpt-dir>/wandb\_uniform\_transition\_diag\_history.jsonl}.
\]
Rows are appended only when
\[
\texttt{train/uniform\_transition\_diag\_active}=1.
\]
This JSONL file is the most convenient source for post-processing the table.

\section{Important Notes}
\begin{itemize}
\item The diagnostic uses the online predictor and current online backbone
hidden states from the same train step.
\item It is an auxiliary measurement only. It is returned through the
\texttt{has\_aux} path and does not alter \(\mathcal{L}_{\mathrm{total}}\).
\item Pair coverage is exhaustive and uniform over the set
\(\{(a,b):1\leq a<b\leq L\}\), so the result does not depend on the current
training pair mode.
\item Timestep coverage is empirical from the current batch; therefore the flag
requires uniform timestep sampling to match the diagnostic table claim.
\end{itemize}

\end{document}
